Hlavním cílem tohoto výzkumného úkolu bylo navrhnout agenta umělé inteligence pro 3D akční hry za pomoci algoritmů strojového učení, díky nimž by se AI dokázala plně přizpůsobit jakékoli situaci, ostatním AI nebo také hráčům.

Nejdříve jsem popsal základní historii umělé inteligence ve hrách a jaké techniky/algoritmy se k tomu nejčastěji používají. Zaměřil jsem se hlavně na AI ve 3D akčních hrách a specificky na hry v tomto žánru, které mají mnohem inteligentnější umělou inteligenci oproti průměrným hrám.

V další části jsem vypsal zjednodušené základy strojového učení, které vedou k popisu všech typů strojového učení, přičemž u každého typu učení byly také vylíčeny některé algoritmy.

Ve třetím úseku jsem popsal návrh samotného autonomního agenta, kde popisuji typ vizualizace, který bych použil, a samotný popis i prostředí, kde by se agent využil. Také zde popisuji a porovnávám dva algoritmy, které by se použily ke sestrojení dvou odlišných částí agenta.

V posledním segmentu jsem vysvětlil metodiku a principy samotného hodnocení agenta, pomocí kterého by se upravovaly vnitřní hodnoty obou algoritmů popsaných v předchozí části.

Tento návrh agenta a prostředí by byl realizován v diplomové práci, kde by byla vysvětlena i samotná konstrukce a zda je tato domněnka vůbec možná sestrojit.